% Metode
\section{Metode}

\subsection{Desain Penelitian}
Paper ini memakai desain deskriptif-korelasional noneksperimental dan tidak mengumpulkan data melalui instrumen primer. Data yang dipakai bersifat sekunder, yaitu skor yang dibagikan secara sukarela oleh peserta UTBK-SNBT 2026 di media sosial X. Estimasi korelasi antarsubtes dihitung dari sampel teramati. Statistik deskriptif sampel, yaitu rerata dan simpangan baku ($SD$), kemudian dibandingkan dengan parameter nasional tidak resmi (belum terverifikasi) untuk menilai sejauh mana sampel mencerminkan populasi sasaran.

\subsection{Sumber Data}
Parameter nasional tidak resmi (belum terverifikasi) yang dipakai sebagai pembanding (Tabel~\ref{tab:deskriptif}) tidak bersumber langsung dari kanal resmi SNPMB. Tabel itu diperoleh dari unggahan di platform X sekitar 8 jam setelah pengumuman skor peserta SNBT 2026. Menurut klarifikasi akun @\textcite{kevins2026}, yang narasi unggahannya mengindikasikan keterlibatan internal kepanitiaan SNPMB walaupun tidak dicantumkan pada profil akun, data itu semula disiapkan untuk materi konferensi pers, tetapi batal ditayangkan atas keputusan rapat internal tanpa penjelasan publik.

Triangulasi independen dilakukan untuk menguji keaslian data itu dengan membandingkan sejumlah nilai ekstrem pada tabel terhadap skor yang diunggah secara mandiri oleh peserta. Hasilnya menunjukkan kesesuaian pada beberapa sampel, seperti skor PK tertinggi sebesar 1.173,55 pada peserta yang diterima di STEI-K ITB dan skor PM tertinggi sebesar 997,70 pada peserta di FTI-RI ITB. Presisi angka yang tinggi dan jeda waktu kemunculan data yang relatif singkat memperkuat dugaan bahwa data itu asli. Namun, triangulasi ini bersifat parsial. Pengujian hanya mencakup dua nilai ekstrem, sedangkan rerata, simpangan baku, dan kuartil belum teruji. Nilai $M = 500$ yang seragam pada seluruh subtes mengikuti definisi skala sehingga tidak informatif sebagai bukti keaslian. Atas dasar itu, tabel ini tetap disebut parameter nasional tidak resmi (belum terverifikasi), dan status perolehannya dicatat secara terbuka pada Bagian Keterbatasan Penelitian.

Berbeda dengan tabel itu, jumlah peserta hadir pada UTBK-SNBT 2026 ($N_{\text{nasional}} = 846.518$, dari 871.496 pendaftar) bersumber dari paparan resmi konferensi pers hasil SNBT 2026 yang dipublikasikan \textcite{snpmb2026konferensi}. Angka ini hanya dipakai untuk menghitung peringkat nasional dan jumlah peserta yang diharapkan berada di luar rentang skor.

Skor pada kedelapan komponen UTBK-SNBT 2026 dihimpun dari unggahan publik di platform X, terutama melalui akun @sbmptnfess dan @utbkcak. Pada akun @sbmptnfess, data dijaring dari utas khusus perguruan tinggi negeri (PTN) tertentu, seperti Institut Teknologi Bandung (ITB), Universitas Brawijaya (UB), Universitas Gadjah Mada (UGM), dan Universitas Indonesia (UI). Data juga dijaring dari utas terbuka yang mengundang pembagian nilai tanpa membatasi status kelulusan. Seluruh data yang dianalisis berupa tangkapan layar sertifikat resmi hasil ujian yang diunggah langsung oleh peserta.

\subsection{Prosedur Pengumpulan Data}
Pengumpulan data berlangsung dalam tiga langkah yang mencakup 19 klaster utas. Pertama, penyisiran utas bertema PTN tertentu di akun @sbmptnfess menjaring $N = 1.459$ entri. Kedua, perluasan ke utas umum pada akun @sbmptnfess menambah jumlah entri menjadi $N = 1.709$. Ketiga, penambahan dari akun @utbkcak menggenapkan total sampel menjadi $N = 1.886$ setelah prosedur pembersihan (lihat Subbagian Pembersihan dan Validasi Data). Tiap penambahan klaster disebut tahap (Tahap 1 sampai 19) pada Bagian Hasil. Bagian itu juga menyajikan tren estimasi korelasi, yang berfluktuasi lalu melandai (\textit{plateau}) seiring penambahan tahap.

\subsection{Variabel Penelitian}
Paper ini menganalisis kedelapan skor subtes UTBK-SNBT 2026, yaitu PU, PPU, PBM, PK, LBI Saintek, LBI Soshum, LBE, dan PM. Skor komposit dihitung sebagai rerata aritmetis kedelapan skor subtes, $Z = \frac{1}{8}\sum_{k=1}^8 X_k$, untuk menyusun matriks translasi persentil standar, tabel yang menerjemahkan titik persentil standar ($M \pm k\,SD$) menjadi skor (lihat Bagian Hasil). LBI Saintek dan LBI Soshum diperlakukan sebagai dua komponen mandiri dengan bobot setara terhadap enam subtes lain, tanpa digabung lebih dahulu menjadi satu nilai.

LBI diujikan sebagai satu tes yang sama bagi seluruh peserta \autocite{detik2026lbi}, sedangkan pada tahun-tahun sebelumnya LBI dilaporkan sebagai satu nilai. Sejauh penelusuran penulis, prosedur Panitia SNPMB menurunkan dua skor dari satu tes itu belum dijelaskan secara terbuka. Dengan demikian, LBI Saintek dan LBI Soshum merupakan dua skor dari satu tes pada peserta yang sama dan tidak independen secara desain. Pada skor komposit, LBI memperoleh bobot $\frac{2}{8}$, dua kali bobot subtes lain ($\frac{1}{8}$). Penulis mempertahankan spesifikasi ini agar sesuai dengan format sertifikat, dan hasil analisis sensitivitasnya dilaporkan pada Bagian Hasil.

\subsection{Pembersihan dan Validasi Data}
Sebelum korelasi diestimasi, data awal disaring dengan tiga kriteria eksklusi. Kriteria pertama adalah \textbf{duplikasi data}, yaitu entri dengan kombinasi nilai identik pada kedelapan subtes, yang dideteksi melalui pencocokan nilai persis antarentri pada skrip pembersihan data. Kriteria kedua adalah \textbf{nilai di luar rentang}, yaitu entri dengan skor yang melampaui batas minimum atau maksimum nasional tidak resmi pada subtes terkait, sebagai indikasi kesalahan input. Kriteria ketiga adalah \textbf{indikasi manipulasi visual atau rekayasa kecerdasan buatan}, yang mencakup (a) keberadaan \textit{watermark} generator kecerdasan buatan pada gambar sertifikat, (b) nilai bulat sempurna, seperti skor tepat 1.000,00, yang tidak mungkin muncul pada penskalaan \textit{Item Response Theory} (IRT), dan (c) ketidakkonsistenan tata letak digit dengan pola sebaran alami sertifikat resmi. Metode IRT menaksir skor dari pola jawaban pada tiap butir soal. Alasannya, skor maksimum nasional PM hanya 997,70, dan PK tidak memiliki riwayat nilai bulat.

Kriteria duplikasi hanya berlaku apabila kedelapan nilai identik sekaligus, sehingga nilai yang sama pada satu subtes saja tidak dianggap duplikasi. Pada LBI Saintek, misalnya, terdapat 1.396 nilai berbeda di antara 1.886 entri (nilai 752,30 muncul 19 kali), tetapi entri-entri itu memiliki profil skor yang berbeda pada tujuh subtes lainnya. Penyebab seringnya nilai kembar pada subtes ini belum ditelusuri.

\subsection{Teknik Analisis Data dan Keterbukaan Replikasi}
Seluruh analisis, estimasi statistik, dan visualisasi dijalankan secara terprogram dengan skrip komputasi terbuka berbahasa pemrograman R versi 4.6.1. Korelasi Pearson teramati dihitung untuk 28 pasangan unik antarsubtes. Rerata korelasi \textit{off-diagonal} (rerata koefisien di luar diagonal matriks korelasi) dihitung dengan membagi jumlah koefisien itu dengan 56, banyak sel di luar diagonal pada matriks $8 \times 8$ ($8^2 - 8 = 56$; Tabel~\ref{tab:korelasi}). Bias seleksi dievaluasi dengan membandingkan rerata dan $SD$ sampel teramati terhadap parameter nasional tidak resmi (Tabel~\ref{tab:deskriptif}).

Penulis memilih formula Thorndike Case III \autocite{thorndike1949,sackett2000} untuk mengoreksi atenuasi (pelemahan korelasi) akibat \textit{restriction of range}. Formula ini merupakan derivasi operasional dari teorema seleksi multivariat Lawley--Aitken \autocite{aitken1934,lawley1943}. Pilihan itu didasarkan pada mekanisme bias sukarela di media sosial yang terjadi langsung pada skor komposit $Z$, sedangkan Case II mengasumsikan seleksi langsung pada salah satu variabel yang dikorelasikan. Akibatnya, korelasi tiap pasangan subtes $X_i$ dan $X_j$ melemah secara tidak langsung (\textit{incidental restriction}), karena seleksi pada $Z$ menjalar ke tiap subtes melalui kovarians keduanya dengan $Z$. Estimasi korelasi populasi terkoreksi ($R_{ij}$) dihitung melalui persamaan:
\begin{equation}
R_{ij} = \frac{r_{ij} + r_{iZ} r_{jZ} \left(\frac{S_Z^2}{s_Z^2} - 1\right)}{\sqrt{\left[1 + r_{iZ}^2 \left(\frac{S_Z^2}{s_Z^2} - 1\right)\right] \left[1 + r_{jZ}^2 \left(\frac{S_Z^2}{s_Z^2} - 1\right)\right]}}
\end{equation}
dengan $r_{ij}$ sebagai korelasi \textit{raw sample}, $r_{iZ}$ dan $r_{jZ}$ sebagai korelasi masing-masing subtes dengan $Z$, $s_Z \approx 53,84$ sebagai $SD$ skor komposit sampel teramati, $S_Z$ sebagai $SD$ skor komposit populasi, dan $\lambda^2 = \frac{S_Z^2}{s_Z^2} - 1$ sebagai multiplikator seleksi (kenaikan relatif varians skor komposit populasi terhadap sampel teramati).

Karena $S_Z$ dan $R_{ij}$ saling bergantung (ketergantungan sirkular), $S_Z$ diestimasi melalui iterasi titik tetap (\textit{fixed-point iteration}), prosedur yang mengulang penghitungan sampai nilainya berhenti berubah. Prosedur ini memakai matriks kovarians heterogen penuh:
\begin{equation}
S_{Z}^{(k+1)} = \sqrt{\frac{1}{64}\left( \sum_{i=1}^8 \sigma_i^2 + \sum_{i=1}^8 \sum_{j \neq i}^8 \sigma_i \sigma_j R_{ij}^{(k)} \right)}
\end{equation}
dengan $\sigma_{\text{PU}} = 71,52$ dan $\sigma_{\text{lain}} = 100,00$ dari parameter nasional tidak resmi (Tabel~\ref{tab:deskriptif}). Iterasi dijalankan pada skrip R dengan toleransi presisi tinggi ($|\Delta S_Z| \le 10^{-12}$), berawal dari $s_Z$, dan mencapai konvergensi stabil pada siklus ke-43.

Agar analisis dapat direproduksi secara komputasi (\textit{computational reproducibility}) dan diperiksa secara terbuka, seluruh kode analisis R, berkas data replikasi tanpa identitas peserta, kamus data (\textit{codebook}), dan skrip pengujian otomatis (\textit{unit tests}) disimpan secara terbuka di repositori riset publik.

\subsection{Analisis Faktor Eksploratori dan Konfirmatori}
Guna menguji apakah kedelapan skor subtes membentuk \textit{g-factor} tunggal atau struktur berjenjang dengan faktor kelompok yang lebih spesifik (lihat Tinjauan Pustaka), penulis menjalankan \textit{exploratory factor analysis} (EFA), yang mencari struktur faktor langsung dari data, dan CFA. Keduanya memakai paket \texttt{psych} \autocite{revelle2024psych} dan \texttt{lavaan} \autocite{rosseel2012lavaan} pada R. \textit{Horn's Parallel Analysis} \autocite{horn1965} dijalankan sebagai diagnostik pendukung. Prosedur ini membandingkan nilai eigen (besar varians yang dijelaskan tiap faktor) matriks korelasi sampel teramati dengan nilai eigen matriks acak hasil simulasi berukuran setara. Penulis menjalankannya berbasis faktor, dengan metode ekstraksi \textit{minres} dan 100 iterasi simulasi. Penulis menetapkan dua faktor pada EFA secara apriori, mengikuti pembedaan domain Literasi dan Numerasi pada kerangka teori (lihat Tinjauan Pustaka), dan tidak semata-mata mengikuti jumlah faktor yang disarankan \textit{Parallel Analysis} (hasilnya dilaporkan pada Bagian Hasil).

Dua model CFA dibandingkan. Model 1 memuat \textit{g-factor} tunggal yang membebani kedelapan subtes dengan muatan yang diestimasi bebas, sedangkan Model 2 memuat dua faktor berkorelasi, yaitu Literasi dan Numerasi, sesuai hasil ekstraksi EFA. Kedua model diestimasi dengan \textit{robust maximum likelihood} (estimator MLM, uji Satorra--Bentler) karena data mengindikasikan sebaran miring dan menyimpang dari normal (lihat Bagian Hasil, Subbagian Evaluasi Asumsi Distribusi). Kecocokan model (kesesuaian antara korelasi yang diturunkan model dan korelasi teramati) dievaluasi melalui $\chi^2_{\text{SB}}$, \textit{Comparative Fit Index} (CFI), \textit{Tucker--Lewis Index} (TLI), \textit{Root Mean Square Error of Approximation} (RMSEA) beserta selang kepercayaan 90\%, \textit{Standardized Root Mean Square Residual} (SRMR), dan \textit{Akaike Information Criterion} (AIC), berdasarkan ambang batas \textcite{hubentler1999}. Keunggulan Model 2 terhadap Model 1 diuji melalui selisih $\chi^2$ berskala (\textit{scaled chi-square difference test}).

\subsection{Analisis Sensitivitas dan Ketidakpastian}
Analisis berikut bersifat tambahan dan tidak mengubah spesifikasi utama, yaitu delapan subtes, CFA dua-faktor apriori, dan $SD$ nasional tidak resmi tanpa penyesuaian.

Pertama, peran LBI sebagai satu tes dengan dua skor diperiksa melalui tiga analisis. Analisis pertama adalah EFA tiga-faktor (ekstraksi \textit{minres}, rotasi \textit{oblimin} yang mengizinkan faktor berkorelasi) untuk menafsirkan faktor ketiga. Analisis kedua adalah CFA dua-faktor dengan kovarians residual (keterkaitan antara LBI Saintek dan LBI Soshum di luar kedua faktor), yang dibandingkan dengan model utama melalui uji selisih $\chi^2$ Satorra--Bentler. Analisis ketiga adalah koreksi Case III yang diulang dengan salah satu atau kedua skor LBI dikeluarkan dari skor komposit. Rerata kedua skor LBI tidak dipakai sebagai satu variabel karena $SD$ nasionalnya tidak tersedia.

Kedua, ketidakpastian pengambilan sampel (variasi estimasi akibat sampel yang ditarik berbeda) diperkirakan dengan \textit{bootstrap} nonparametrik, yang menarik ulang entri secara acak dari sampel teramati. Sebanyak 2.000 sampel \textit{bootstrap} ditarik dengan pengembalian dari seluruh entri sampel teramati (benih acak 2026). Pada tiap sampel \textit{bootstrap}, korelasi \textit{raw sample}, korelasi subtes-komposit, dan $SD$ skor komposit dihitung ulang sebelum iterasi titik tetap dijalankan sampai konvergen. Interval persentil 95\% dilaporkan untuk $S_Z$, $\lambda^2$, dan $\bar{R}$. \textit{Bootstrap} hanya menangkap ketidakpastian pengambilan sampel.

Ketiga, ketidakakuratan $SD$ nasional diuji dengan mengulang estimasi utama pada seluruh entri sampel teramati setelah seluruh $SD$ nasional dikalikan $0,90$; $0,95$; $1,05$; atau $1,10$. Pengulangan juga dilakukan setelah $SD_{\text{PU}}$ diganti $100$ dengan $SD$ subtes lain tetap.

Keempat, asumsi normalitas yang mendasari translasi persentil diperiksa dalam dua langkah. Langkah pertama adalah diagnostik deskriptif yang dihitung dari tujuh statistik nasional tidak resmi tiap subtes. Diagnostik ini mencakup skor-$z$ minimum dan maksimum, jumlah peserta yang diharapkan berada di luar rentang tersebut pada distribusi normal dengan 846.518 peserta hadir, kemiringan Bowley $\frac{Q_3 + Q_1 - 2\,\text{Med}}{Q_3 - Q_1}$ dengan $Q_1$, Med, dan $Q_3$ sebagai kuartil pertama, median, dan kuartil ketiga, selisih rerata dan median, dan rasio rentang antarkuartil terhadap $SD$ ($\frac{IQR}{SD}$), yang bernilai $1,349$ pada distribusi normal.

Langkah kedua adalah simulasi. Karena data mikro nasional (skor tiap peserta) tidak tersedia, transformasi normalisasi seperti Box-Cox tidak dapat diterapkan. Sebagai gantinya, distribusi marginal tiap subtes (sebaran skor satu subtes secara terpisah) dihampiri distribusi Johnson SB \autocite{johnson1949}. Distribusi berbatas dengan empat parameter ini dapat meniru ekor kiri terpotong dan ekor kanan panjang. Parameter dicocokkan dengan kuadrat terkecil berbobot terhadap tujuh statistik nasional tidak resmi. Skor minimum dan maksimum masing-masing diletakkan pada peluang $\frac{0,5}{N}$ dan $1 - \frac{0,5}{N}$ sebagai hampiran nilai ekstrem pada sampel besar, dan pencarian dimulai dari 40 titik awal acak (benih acak 2026). Skor komposit kemudian disimulasikan untuk 1.000.000 peserta virtual dengan \textit{copula} Gaussian, model yang menghubungkan distribusi marginal melalui satu korelasi laten seragam pada skala normal baku. Simulasi A memakai korelasi laten sebesar $\bar{R}$ hasil Thorndike Case III. Simulasi B memakai korelasi laten yang dikalibrasi agar $SD$ skor komposit simulasi sama dengan $S_Z$, sehingga pengaruh bentuk distribusi terpisah dari pengaruh skala. Skor komposit simulasi dipusatkan ke rerata $500$ karena hampiran marginal menyisakan galat rerata beberapa poin. Hasil simulasi dibandingkan dengan versi normal melalui dua ukuran, yaitu peluang ekor di luar titik $M \pm k\,S_Z$ dan skor komposit simulasi pada persentil yang sama dengan persentil normal di titik tersebut.

Kelima, \textit{Spearman's Law of Diminishing Returns} (SLODR), atau diferensiasi kemampuan, diperiksa melalui model satu faktor. SLODR menyatakan bahwa korelasi antartes lebih rendah pada kemampuan tinggi dan lebih tinggi pada kemampuan rendah. Muatan \textit{g-factor} pada model itu menurun linear terhadap tingkat kemampuan, $m(g) = m_0 - b\,g$ dengan $g \sim N(0,1)$ dan $b \in \{0; 0,01; 0,02; 0,03; 0,04\}$ ($b = 0$ berarti korelasi konstan). Tiap skor laten diubah menjadi normal baku melalui peringkat agar distribusi marginal tidak berubah. Simulasi ini terdiri atas dua bagian. Pada bagian pertama, marginal Johnson SB digunakan dan $m_0$ dikalibrasi agar $SD$ skor komposit sama dengan $S_Z$, lalu peluang ekor dan skor komposit dibandingkan antarskenario dan dengan normal. Pada bagian kedua, populasi Gaussian dengan $SD$ nasional tidak resmi diseleksi lewat skor komposit dengan probabilitas seleksi logistik, yang parameternya dicocokkan ke rerata dan $SD$ skor komposit sampel teramati. Nilai $m_0$ dikalibrasi agar rerata korelasi \textit{raw sample} pada sampel simulasi sama dengan sampel teramati ($b \in \{0; 0,02; 0,04\}$). Prosedur Case III yang sama kemudian dijalankan pada sampel simulasi dan dibandingkan dengan nilai sebenarnya pada populasi simulasi. Simulasi SLODR memakai $300.000$ peserta virtual dan benih acak 2026.

\subsection{Pertimbangan Etis}
Seluruh data yang dianalisis bersumber dari unggahan publik pada platform X yang dapat diakses secara terbuka tanpa autentikasi khusus. Analisis dan pelaporan hasil dilakukan sepenuhnya pada tingkat agregat, tanpa mencantumkan identitas pribadi peserta (nama akun maupun identitas asli) di bagian mana pun dalam paper ini.